\documentclass[12pt]{extarticle}
% Shared preamble for the Week 1 student pages.
\usepackage[letterpaper,left=0.55in,right=0.55in,top=0.85in,bottom=0.8in,
            headheight=18pt,headsep=14pt,footskip=28pt]{geometry}
\usepackage[T1]{fontenc}
\usepackage[default]{andika}
\usepackage{tikz}
\usetikzlibrary{arrows.meta}
\usepackage{fancyhdr}
\usepackage{xcolor}

\definecolor{gridgray}{gray}{0.60}
\definecolor{chainfill}{gray}{0.78}
\definecolor{holefill}{gray}{0.08}
\definecolor{blockG}{HTML}{3DAE4B}
\definecolor{blockB}{HTML}{3F6FD8}
\definecolor{blockR}{HTML}{E2463B}
\definecolor{blockY}{HTML}{F6D12E}
\definecolor{blockP}{HTML}{8A4FC0}
\colorlet{fillG}{blockG!25}
\colorlet{fillB}{blockB!18}
\colorlet{fillR}{blockR!25}
\colorlet{fillY}{blockY!35}
\colorlet{fillP}{blockP!25}
\tikzset{writeline/.style={black!45,line width=0.5pt}}

\pagestyle{fancy}
\fancyhf{}
\renewcommand{\headrulewidth}{0.4pt}
\renewcommand{\footrulewidth}{0pt}
\lfoot{Bellingham Math Circle / Week 1}
\rfoot{\thepage}

\setlength{\parindent}{0pt}
\setlength{\parskip}{0pt}
\raggedright

% "Problem N:" label followed by ordinary sentences.
\newcommand{\problem}[1]{\textbf{Problem #1:}\ }
\newcommand{\fig}[1]{\par\noindent\input{figs/#1}\par}

% A block picture followed by its name, for the "which blocks" line of a problem.
\newcommand{\blk}[2]{\mbox{\input{figs/icon_#1}\ \ #2}}
\newcommand{\blkgap}{\hspace{0.45in}}

\lhead{Week 1 / Tiling / Grades 4--5}
\linespread{1.1}

\begin{document}

For every problem: use blue rhombuses only. To cover a board, fill all of it with
no gaps, no overlaps, and nothing sticking out. Two coverings are different if
their drawings are different, even when one is a turned or mirror copy of the
other.

\bigskip
\problem{1}Find all the different ways to cover board A, and draw each one in the
small pictures beside the board. Do the same for board B and board C. Then
imagine a board shaped like these whose top edge is 10 small edges long. How
many ways are there to cover it? Explain how you know.

\vspace{0.3in}
\fig{g45_p1}

\newpage
Board with a top edge 10 small edges long: \underline{\hspace{0.8in}} ways
\fig{g45_lines_p1}

\vspace{0.4in}
\problem{2}Find all the different ways to cover board D, and draw each one.
Explain how you know you have found them all.

\vspace{0.3in}
\fig{g45_p2}
\fig{g45_lines_p2}

\newpage
\problem{3}Cover board E to match picture S. Then change it into picture T, one
move at a time. In one move, you find three rhombuses that together make the
shape of a yellow hexagon, lift them out, and put them back the other way. What is the
fewest moves that takes S to T? Do the same from S to U, and from U to T.

\vspace{0.3in}
\fig{g45_p3}

\vspace{0.3in}
S to T: \underline{\hspace{0.7in}} moves \hfill
S to U: \underline{\hspace{0.7in}} moves \hfill
U to T: \underline{\hspace{0.7in}} moves

\newpage
\problem{4}In a covering of a board, start at one small edge on the top side.
Go down through the rhombus below it, out through that rhombus's bottom edge, and
into the next rhombus. Keep going until you reach the bottom side. Shade this
chain of rhombuses, and write a letter for each rhombus in it, from top to
bottom: R if it leans right, L if it leans left. The picture shows one chain in
a covering of board D.

\vspace{0.2in}
\fig{g45_p4_rl}

\vspace{0.2in}
Shade both chains in each covering of board D that you drew in Problem 2, and
write their letters. Then, for each pair below, draw the covering of board E
whose left chain and right chain have these letters, or explain why there is
none.

\vspace{0.3in}
\fig{g45_p4_pairs}
\fig{g45_lines_p4}

\newpage
\problem{5}How many different ways are there to cover board E? Draw them.
Explain how you know you have found them all.

\vspace{0.3in}
\fig{g45_p5}

\vspace{0.3in}
Number of ways: \underline{\hspace{0.8in}}
\fig{g45_lines_p5}

\newpage
\problem{6}Cover board E any way you like, and write the letters of its two
chains. Make one move, as in Problem 3, and write the letters again. Try this
with many different moves. What does one move do to the letters?

\fig{g45_lines_p6}

\vspace{0.4in}
\problem{7}Explain why it takes at least 8 moves to get from S to T. Then start
with any covering of board E, make some moves, and finish with the same covering
you started with. Can the number of moves be odd? Explain.

\fig{g45_lines_p7}

\end{document}
